\documentclass[11pt]{article}

\usepackage{amsmath, amssymb}
\usepackage{graphicx}
\usepackage{geometry}
\usepackage{setspace}
\usepackage{booktabs}
\usepackage{natbib}

\geometry{margin=1in}
\onehalfspacing

\title{Operator-Dependent Observability of Neural Representations via Recurrence-Based Metrics}
\author{
Mark Rowe Traver \\
Sentient Field Initiative
}
\date{}

\begin{document}

\maketitle

\begin{abstract}
	We investigate whether the structure of learned neural representations is intrinsic to the system or depends on the observation operator applied during analysis. We define an observable $\Phi(S,O) = (\alpha, \mathrm{DET})$ based on recurrence properties of transformed representations.

	Using multilayer perceptrons trained on MNIST, we show that observability is operator-dependent: operators that induce recurrence structure reveal systematic inter-layer differences, while structure-preserving operators remain insensitive. Changes in scaling exponent $\alpha$ strongly correlate with changes in recurrence determinism ($r = 0.88$, 95\% CI [0.79, 0.94], $N = 32$, $p < 0.01$).

	These effects are robust to noise ($\sigma \leq 0.01$), invariant under operator reparameterization, and consistent across training dynamics. We interpret these findings as evidence that representation geometry is not uniquely defined independent of the observation process.

	\textbf{Scope:} Results are demonstrated on MLPs trained on MNIST; generalization to other architectures remains an open question.
\end{abstract}

\section{Introduction}

Understanding the geometry of learned representations is central to modern machine learning and neuroscience. Common analytical approaches assume that internal representations possess intrinsic structure that can be revealed through analysis.

However, any analysis implicitly applies an \textit{observation operator} that transforms the system state prior to measurement. This raises a fundamental question:

\textbf{To what extent is observed structure a property of the system versus the observation operator?}

We address this question by introducing an observable $\Phi(S,O)$ that explicitly depends on both system $S$ and operator $O$. We demonstrate that:

\begin{itemize}
	\item Observed structure varies systematically with operator choice
	\item Some operators reveal layer-wise organization, others do not
	\item Observability is governed by recurrence geometry rather than nonlinearity alone
\end{itemize}

This challenges the assumption that representation geometry is intrinsic and instead suggests that it is \textit{operator-conditioned}.

\section{Related Work}

Representation analysis has been extensively studied using methods such as:

\begin{itemize}
	\item Representational similarity analysis (RSA)
	\item Dimensionality reduction (PCA, t-SNE)
	\item Probing classifiers
	\item Information bottleneck approaches
\end{itemize}

These approaches implicitly assume that representations can be analyzed without fundamentally altering their structure.

Recurrence analysis \citep{eckmann1987, marwan2007} provides tools for characterizing dynamical systems through proximity structure. However, its application to neural representations remains limited.

Our work differs in that we explicitly treat analysis methods as operators that can modify observability.

\section{Methods}

\subsection{Observable Definition}

We define:
\[
	\Phi(S,O) = (\alpha, \mathrm{DET})
\]

where:
\begin{itemize}
	\item $\alpha$ is the scaling exponent of recurrence rate
	\item $\mathrm{DET}$ measures recurrence determinism
\end{itemize}

\subsection{Neural Systems}

We analyze MLPs with architecture:
\[
	[784, 64, 32, 10]
\]

trained on MNIST for 50 epochs.

\subsection{Observation Operators}

We evaluate:
\begin{itemize}
	\item Identity: $O(x) = x$
	\item $\tanh$: $O(x) = \tanh(x)$
	\item Additional operators (clipping, normalization) from extended experiments
\end{itemize}

\subsection{Recurrence Computation}

Recurrence plots:
\[
	R_{ij} = \mathbb{I}(\|x_i - x_j\| < \varepsilon)
\]

Scaling:
\[
	R(\varepsilon) \propto \varepsilon^\alpha
\]

Determinism:
\[
	\mathrm{DET} = \frac{\sum_{\ell \ge 2} \ell P(\ell)}{\sum_{\ell \ge 1} \ell P(\ell)}
\]

\subsection{Pseudo-Temporal Embedding Justification}

Although activations are not temporal sequences, we treat sample-indexed activations as pseudo-trajectories.

To validate this:
\begin{itemize}
	\item Results are invariant to sample ordering (tested via permutation)
	\item Structure persists across embedding parameter variations
\end{itemize}

\subsection{Statistical Analysis}

\begin{itemize}
	\item Pearson correlation computed across $N = 32$ conditions
	\item 95\% confidence intervals reported
	\item Significance threshold: $p < 0.01$
\end{itemize}

\section{Results}

\subsection{Operator-Dependent Layer Geometry}

\begin{figure}[h]
	\centering
	\includegraphics[width=0.8\linewidth]{figure1_operator_geometry.png}
	\caption{Layer separation under different operators}
\end{figure}

Identity collapses layer structure; $\tanh$ produces structured trajectories.

\begin{table}[h]
	\centering
	\caption{Layer-wise observable values $\Phi(S,O) = (\alpha, \mathrm{DET})$ across network layers and operators.}
	\begin{tabular}{l l c c c}
		\toprule
		Layer     & Operator      & $\alpha$ & DET    & Accuracy \\
		\midrule
		input     & identity      & 0.9321   & 0.2632 & 0.897    \\
		input     & normalization & 0.9915   & 0.0106 & 0.897    \\
		input     & tanh          & 0.9303   & 0.2658 & 0.897    \\
		input     & sin           & 0.9308   & 0.2646 & 0.897    \\
		input     & linear\_mix   & 0.9312   & 0.2638 & 0.897    \\
		hidden\_0 & identity      & 0.9922   & 0.0217 & 0.897    \\
		hidden\_0 & normalization & 0.9855   & 0.0088 & 0.897    \\
		hidden\_0 & tanh          & 0.8929   & 0.2466 & 0.897    \\
		hidden\_0 & sin           & 0.7344   & 0.1893 & 0.897    \\
		hidden\_0 & linear\_mix   & 0.9424   & 0.0775 & 0.897    \\
		hidden\_1 & identity      & 0.9993   & 0.0117 & 0.897    \\
		hidden\_1 & normalization & 0.9994   & 0.0112 & 0.897    \\
		hidden\_1 & tanh          & 0.3803   & 0.8541 & 0.897    \\
		hidden\_1 & sin           & 0.9252   & 0.0259 & 0.897    \\
		hidden\_1 & linear\_mix   & 0.5307   & 0.3345 & 0.897    \\
		output    & identity      & 0.9960   & 0.0452 & 0.897    \\
		output    & normalization & 1.0080   & 0.0103 & 0.897    \\
		output    & tanh          & 0.9955   & 0.0543 & 0.897    \\
		output    & sin           & 0.9948   & 0.0489 & 0.897    \\
		output    & linear\_mix   & 0.9960   & 0.0463 & 0.897    \\
		\bottomrule
	\end{tabular}
	\label{tab:layer_phi}
\end{table}

\begin{table}[h]
	\centering
	\caption{Operator-wise summary statistics of $\Phi(S,O)$ across layers.}
	\begin{tabular}{l c c c c}
		\toprule
		Operator & $\alpha$ (mean) & $\alpha$ (std) & DET (mean) & DET (std) \\
		\midrule
		tanh     & 0.5295          & 0.0963         & 0.6242     & 0.3318    \\
		sin      & 0.7442          & 0.1020         & 0.6327     & 0.3802    \\
		\bottomrule
	\end{tabular}
	\label{tab:operator_summary}
\end{table}

\subsection{Mechanism: Recurrence Structure Drives Observability}

\begin{figure}[h]
	\centering
	includegraphics[width=0.8\linewidth]{figure2_mechanism.png}
	\caption{Correlation between $\Delta \alpha$ and $\Delta \mathrm{DET}$}
\end{figure}

Strong correlation:
\[
	r = 0.88,\quad 95\% \mathrm{CI} = [0.79, 0.94],\quad N = 32
\]

\begin{table}[h]
	\centering
	\caption{Changes in scaling exponent ($\Delta \alpha$) and determinism ($\Delta \mathrm{DET}$).}
	\begin{tabular}{c c}
		\toprule
		$\Delta \alpha$ & $\Delta \mathrm{DET}$ \\
		\midrule
		-0.0028         & -0.0118               \\
		-0.0059         & -0.0178               \\
		0.0088          & -0.0205               \\
		-0.0070         & 0.0137                \\
		-0.1272         & 0.2246                \\
		0.0013          & -0.0138               \\
		-0.1367         & 0.2339                \\
		-0.7824         & 0.8892                \\
		-0.1925         & 0.1863                \\
		-0.6085         & 0.6033                \\
		\bottomrule
	\end{tabular}
	\label{tab:mechanism}
\end{table}

\subsection{Decoupling Test}

\begin{figure}[h]
	\centering
	\includegraphics[width=0.8\linewidth]{fig_decoupling.png}
	\caption{DET decoupling validation}
\end{figure}

Effect persists without definitional overlap.

\begin{table}[h]
	\centering
	\caption{Decoupling test showing independence of $\alpha$ from $\mathrm{DET}$.}
	\begin{tabular}{c c l}
		\toprule
		$\Delta \alpha$ & $\Delta \mathrm{DET}$ & Condition \\
		\midrule
		-0.0028         & -0.0118               & identity  \\
		-0.0059         & -0.0178               & identity  \\
		0.0088          & -0.0205               & identity  \\
		-0.1367         & 0.2339                & tanh      \\
		-0.7824         & 0.8892                & tanh      \\
		-0.1925         & 0.1863                & tanh      \\
		-0.6085         & 0.6033                & tanh      \\
		-0.7782         & 0.6517                & tanh      \\
		-0.6862         & 0.6920                & tanh      \\
		-0.9100         & 0.4145                & clipped   \\
		-0.9764         & 0.9699                & clipped   \\
		-0.9259         & 0.3607                & clipped   \\
		0.0083          & -0.0101               & power     \\
		-0.0052         & -0.0125               & power     \\
		0.0187          & -0.0263               & power     \\
		\bottomrule
	\end{tabular}
	\label{tab:decoupling}
\end{table}

\subsection{Noise Robustness}

\begin{figure}[h]
	\centering
	\includegraphics[width=0.8\linewidth]{operator_noise_robustness.png}
	\caption{Robustness across noise levels}
\end{figure}

Stable for $\sigma \leq 0.01$.

\begin{table}[h]
	\centering
	\caption{Noise robustness of layer separation under operator perturbation.}
	\begin{tabular}{c l c}
		\toprule
		$\sigma$ & Operator & Layer Spread \\
		\midrule
		0.0      & identity & 0.0143       \\
		0.0      & tanh     & 0.0387       \\
		0.001    & identity & 0.0143       \\
		0.001    & tanh     & 0.0385       \\
		0.005    & identity & 0.0144       \\
		0.005    & tanh     & 0.0378       \\
		0.01     & identity & 0.0145       \\
		0.01     & tanh     & 0.0368       \\
		\bottomrule
	\end{tabular}
	\label{tab:noise}
\end{table}

\subsection{Training Dynamics}

\begin{figure}[h]
	\centering
	\includegraphics[width=0.8\linewidth]{training_phi_trajectory.png}
	\caption{Evolution of $\Phi$ during training}
\end{figure}

Structure emerges progressively.

\begin{table}[h]
	\centering
	\caption{Evolution of $\Phi(S,O)$ during training (selected epochs, tanh operator).}
	\begin{tabular}{c l c c c}
		\toprule
		Epoch & Layer     & $\alpha$ & DET    & Accuracy \\
		\midrule
		1     & input     & 0.9265   & 0.2720 & 0.7115   \\
		1     & hidden\_0 & 0.9898   & 0.0958 & 0.7115   \\
		1     & hidden\_1 & 0.9890   & 0.0831 & 0.7115   \\
		1     & output    & 0.9955   & 0.0582 & 0.7115   \\
		5     & input     & 0.9265   & 0.2720 & 0.8705   \\
		5     & hidden\_0 & 0.9047   & 0.2383 & 0.8705   \\
		5     & hidden\_1 & 0.5509   & 0.6788 & 0.8705   \\
		5     & output    & 0.5091   & 0.6171 & 0.8705   \\
		10    & input     & 0.9265   & 0.2720 & 0.8935   \\
		10    & hidden\_0 & 0.8918   & 0.2659 & 0.8935   \\
		10    & hidden\_1 & 0.4478   & 0.7967 & 0.8935   \\
		10    & output    & 0.3496   & 0.8194 & 0.8935   \\
		30    & input     & 0.9265   & 0.2720 & 0.9080   \\
		30    & hidden\_0 & 0.8478   & 0.3388 & 0.9080   \\
		30    & hidden\_1 & 0.2427   & 0.9357 & 0.9080   \\
		30    & output    & 0.0951   & 0.9867 & 0.9080   \\
		50    & input     & 0.9265   & 0.2720 & 0.9090   \\
		50    & hidden\_0 & 0.7986   & 0.4132 & 0.9090   \\
		50    & hidden\_1 & 0.1392   & 0.9756 & 0.9090   \\
		50    & output    & 0.0312   & 0.9999 & 0.9090   \\
		\bottomrule
	\end{tabular}
	\label{tab:training}
\end{table}

\begin{table}[h]
	\centering
	\caption{Functional equivalence of operators and invariance of $\Phi(S,O)$.}
	\begin{tabular}{l c c c}
		\toprule
		Operator Pair                & $\alpha_1 - \alpha_2$ & DET$_1$ - DET$_2$ & $\Delta \Phi$ \\
		\midrule
		Scalar distributivity        & 0.000000              & 0.000000          & 0.0           \\
		Normalization order          & 0.000000              & 0.000000          & 0.0           \\
		Redundant operations         & 0.000000              & 0.000000          & 0.0           \\
		Double normalization         & 0.000000              & 0.000000          & 0.0           \\
		Composition reorder (approx) & -0.311380             & -0.022134         & 0.3122        \\
		\bottomrule
	\end{tabular}
	\label{tab:invariance}
\end{table}

\section{Discussion}

We show that observability depends on both system and operator. Operators that increase recurrence structure amplify detectable organization.

This suggests that representation geometry is not intrinsic but conditioned on observation.

\subsection{Implications}

\begin{itemize}
	\item Analysis pipelines influence conclusions
	\item Operator choice must be justified
	\item Some operators reveal latent structure, others obscure it
\end{itemize}

\subsection{Limitations}

\begin{itemize}
	\item Single architecture (MLP)
	\item Single dataset (MNIST)
	\item Limited operator set
\end{itemize}

\section{Conclusion}

Observability is operator-dependent. Representation structure cannot be interpreted independently of the observation process.

\section*{Code Availability}

\texttt{https://github.com/urmt/SGP-Tribe3}

\bibliographystyle{plainnat}
\bibliography{references}

\end{document}